\documentclass{beamer}
\usepackage{graphicx}
\usepackage{xcolor}
\usepackage{tikz}

% 自定义颜色
\definecolor{purpleblue}{RGB}{80, 80, 180}
\definecolor{lightgray}{RGB}{230, 230, 230}
\definecolor{novelPurple}{RGB}{98, 54, 255} 
\definecolor{darkgray}{RGB}{50, 50, 50}

% 主题设置
\usetheme{Madrid}  
\usecolortheme{dolphin} 

% 设置字体（移除 fontspec，确保 pdfLaTeX 兼容）
\setbeamerfont{title}{size=\Large,series=\bfseries}
\setbeamerfont{frametitle}{size=\LARGE,series=\bfseries}
\setbeamercolor{frametitle}{fg=white, bg=novelPurple}
\setbeamercolor{title}{fg=white, bg=purpleblue}

% 封面信息
\title[SAT Unit 2.1.1 Exponents \& Radicals]{\textbf{SAT Unit 2.1.1\\ Exponents \& Radicals}}
\author{\textbf{Jack Zeng}}
\institute{\textcolor{white}{\textbf{Novel Prep}}}
\date{\today}
\begin{document}

% --- Title Slide ---
\begin{frame}
    \titlepage
    \vfill
    \begin{flushright}
        \textcolor{purpleblue}{\Huge \textbf{Novel Prep}}
    \end{flushright}
\end{frame}


\begin{frame}
\section{Exponent Expression}
    \begin{tikzpicture}[remember picture, overlay]

        \shade[top color=softPurple, bottom color=softGray] 
            (current page.north west) rectangle (current page.south east);

        \fill[white, opacity=0.3] (current page.center) circle (4cm);
        

        \node[align=center] at (current page.center) {
            {\Huge\bfseries\textcolor{purpleblue}{Novel Prep}} \\[0.5cm]
            {\Large\itshape\textcolor{lightText}{Excellence in SAT Prep}}
        };
        \foreach \x in {1,2,3,4,5} {
            \fill[purpleblue, opacity=0.2] (rand*10-5, rand*7-3.5) circle (0.1+\x*0.05);
        }
        ;
    \end{tikzpicture}
\end{frame}


\begin{frame}{Overview}
    \tableofcontents
\end{frame}


% Slide: Title Page Already Included
% Slide: Second Title Design Already Included

\begin{frame}
    \begin{tikzpicture}[remember picture, overlay]

        \shade[top color=softPurple, bottom color=softGray] 
            (current page.north west) rectangle (current page.south east);

        \fill[white, opacity=0.3] (current page.center) circle (4cm);
        

        \node[align=center] at (current page.center) {
            {\Huge\bfseries\textcolor{purpleblue}{Exponent}} \\[0.5cm]
            {\Large\itshape\textcolor{lightText}{Excellence in SAT Prep}}
        };
        \foreach \x in {1,2,3,4,5} {
            \fill[purpleblue, opacity=0.2] (rand*10-5, rand*7-3.5) circle (0.1+\x*0.05);
        }
        ;
    \end{tikzpicture}
\end{frame}

% --- Frame 1: Definition ---
\begin{frame}{Definition of Exponents}
\textbf{Definition:} An exponential expression has the form $x^n$, where:
\begin{itemize}
    \item $x$ is the \textbf{base}
    \item $n$ is the \textbf{exponent} or power
\end{itemize}

\vspace{0.3cm}
This represents the base $x$ being multiplied by itself $n$ times.

\[
x^n = \underbrace{x \cdot x \cdot x \cdots x}_{n \text{ times}}
\]
\end{frame}

% --- Frame 2: Laws of Exponents ---
\begin{frame}{Laws of Exponents (Part 1)}

\begin{block}{Basic Properties}

\begin{itemize}
    \item \textbf{Identity:} \quad $x^1 = x$ \hfill (e.g., $3^1 = 3$)
    \item \textbf{Zero Exponent:} \quad $x^0 = 1$ \hfill (e.g., $3^0 = 1$)
\end{itemize}
\end{block}

\begin{block}{Multiplication and Division Rules}
\vspace{0.2cm}
\begin{itemize}
    \item \textbf{Product Rule:} \quad $x^m \cdot x^n = x^{m+n}$ \hfill (e.g., $3^4 \cdot 3^5 = 3^9$)
    \item \textbf{Quotient Rule:} \quad $\dfrac{x^m}{x^n} = x^{m-n}$ \hfill (e.g., $\dfrac{3^7}{3^3} = 3^4$)
\end{itemize}
\end{block}

\end{frame}

\begin{frame}{Laws of Exponents (Part 2)}

\begin{block}{Power Rules}

\begin{itemize}
    \item \textbf{Power of a Power:} \quad $(x^m)^n = x^{mn}$ \hfill (e.g., $(3^2)^4 = 3^8$)
    \item \textbf{Power of a Product:} \quad $(xy)^m = x^m y^m$ \hfill (e.g., $(2 \cdot 3)^3 = 2^3 \cdot 3^3$)
    \item \textbf{Power of a Quotient:} \quad $\left(\dfrac{x}{y}\right)^m = \dfrac{x^m}{y^m}$ \hfill (e.g., $\left(\dfrac{2}{3}\right)^3 = \dfrac{2^3}{3^3}$)
\end{itemize}
\end{block}

\begin{block}{Negative Exponent}

\begin{itemize}
    \item \textbf{Rule:} \quad $x^{-m} = \dfrac{1}{x^m}$ \hfill (e.g., $3^{-4} = \dfrac{1}{3^4}$)
\end{itemize}
\end{block}

\end{frame}

\begin{frame}{\textbf{Question: Product of Exponential Expressions}}
Which expression represents the product of 
\[
(b^6 c^{-2} d^{-5}) \quad \text{and} \quad (b^8 c^{-3}+c^4 d^5),
\]
where \( b, c, \) and \( d \) are positive?

\vspace{0.5em}
\textbf{Options:}
\begin{enumerate}
    \item[A.] \( b^{14} c^{-5} + c^2 d^{-10} \)
    \item[B.] \( b^{14} c^{-5} + c^2 \)
    \item[C.] \( b^{14} c^{-5} d^{-5} + b^6 c^2 \)
    \item[D.] \( b^{14} c^{-5} d^{-5} + c^2 \)
\end{enumerate}
\end{frame}


\begin{frame}{\textbf{Answer and Explanation}}


\textbf{Explanation:}

We are asked to find the product of:
\[
(b^6 c^{-2} d^{-5}) \quad \text{and} \quad (b^8 c^{-3} + c^4 d^5)
\]

Distribute the first term across both terms in the second expression:
\[
b^6 c^{-2} d^{-5} \cdot b^8 c^{-3} = b^{6+8} c^{-2-3} d^{-5} = b^{14} c^{-5} d^{-5}
\]
\[
b^6 c^{-2} d^{-5} \cdot c^4 d^5 = b^6 c^{2} d^{0} = b^6 c^2
\]

Putting it together:
\[
b^{14} c^{-5} d^{-5} + b^6 c^2
\]

Now check the answer choices. Only choice \textbf{C} has this exact expression.

\textbf{Final Answer:} \textbf{C.} \( b^{14} c^{-5} d^{-5} + b^6 c^2 \)
\end{frame}


\begin{frame}
\section{Radical \& Roots }
    \begin{tikzpicture}[remember picture, overlay]

        \shade[top color=softPurple, bottom color=softGray] 
            (current page.north west) rectangle (current page.south east);

        \fill[white, opacity=0.3] (current page.center) circle (4cm);
        

        \node[align=center] at (current page.center) {
            {\Huge\bfseries\textcolor{purpleblue}{Radical \& Roots}} \\[0.5cm]
            {\Large\itshape\textcolor{lightText}{Excellence in SAT Prep}}
        };
        \foreach \x in {1,2,3,4,5} {
            \fill[purpleblue, opacity=0.2] (rand*10-5, rand*7-3.5) circle (0.1+\x*0.05);
        }
        ;
    \end{tikzpicture}
\end{frame}

% --- Frame 4: Definition of Radicals ---
\begin{frame}{Radicals and Roots}
\textbf{Definition:} A radical is the inverse of an exponent.

\[
\sqrt[n]{x} = x^{\frac{1}{n}}
\]

\begin{itemize}
    \item $\sqrt{x} = x^{\frac{1}{2}}$ (square root)
    \item $\sqrt[3]{x} = x^{\frac{1}{3}}$ (cube root)
    \item $\sqrt[4]{x^3} = x^{\frac{3}{4}}$ (fractional exponent)
\end{itemize}
\end{frame}

% --- Frame 5: Properties of Radicals ---
\begin{frame}{Properties of Radicals}
\begin{itemize}
    \item $\sqrt{a} \cdot \sqrt{b} = \sqrt{ab}$
    \item $\dfrac{\sqrt{a}}{\sqrt{b}} = \sqrt{\dfrac{a}{b}}$
    \item $\sqrt[n]{x^m} = x^{\frac{m}{n}}$
\end{itemize}

\vspace{0.2cm}
\textbf{Important:} All radicals can be rewritten using fractional exponents!
\end{frame}

% --- Frame 6: Rationalizing Denominators ---
\begin{frame}{Rationalizing Denominators}
\textbf{Goal:} Remove square roots from denominators.

\[
\dfrac{1}{\sqrt{2}} \quad \Rightarrow \quad \dfrac{1}{\sqrt{2}} \cdot \dfrac{\sqrt{2}}{\sqrt{2}} = \dfrac{\sqrt{2}}{2}
\]

\vspace{0.3cm}
Use the \textbf{conjugate} if the denominator is a binomial:
\[
\dfrac{1}{\sqrt{2} + 1} \cdot \dfrac{\sqrt{2} - 1}{\sqrt{2} - 1} = \dfrac{\sqrt{2} - 1}{(\sqrt{2})^2 - 1^2} = \dfrac{\sqrt{2} - 1}{1}
\]
\end{frame}

% --- Frame 7: Challenge Problem ---
\begin{frame}{Question}
\textbf{Simplify the expression:}
\[
\left( \dfrac{27x^6 y^3}{81x^2 y^9} \right)^{\frac{1}{3}}
\]

\end{frame}

\begin{frame}{Solution}
    

\textbf{Solution:}
\[
= \left( \dfrac{27}{81} \cdot x^{6-2} \cdot y^{3-9} \right)^{\frac{1}{3}} 
= \left( \dfrac{1}{3} \cdot x^4 \cdot y^{-6} \right)^{\frac{1}{3}} 
= \dfrac{x^{\frac{4}{3}}}{3 \cdot y^2}
\]
\end{frame}

\begin{frame}{\textbf{Question: Rational Exponents and Roots}}
If 
\[
\frac{\sqrt{x^5}}{\sqrt[3]{x^4}} = x^{\frac{a}{b}}
\]
for all positive values of \( x \), what is the value of \( \frac{a}{b} \)?

\end{frame}

\begin{frame}{Solution}
\textbf{Solution:}

We use the property of radicals and exponents:
\[
\sqrt{x^5} = x^{5/2}, \quad \sqrt[3]{x^4} = x^{4/3}
\]
\[
\frac{x^{5/2}}{x^{4/3}} = x^{5/2 - 4/3}
\]

Find a common denominator:
\[
\frac{5}{2} = \frac{15}{6}, \quad \frac{4}{3} = \frac{8}{6}
\]
\[
x^{15/6 - 8/6} = x^{7/6}
\]

Thus,
\[
\frac{a}{b} = \frac{7}{6}
\]

\textbf{Answer:} \( \boxed{\frac{7}{6}} \)
    
\end{frame}


\begin{frame}{\textbf{Question: Exponent and Radical Expression Equivalence}}
For what value of \( x \) is the given expression equivalent to \( (70n)^{30x} \), where \( n > 1 \)?
\[
\sqrt[5]{70n \left( \sqrt[6]{70n} \right)^2}
\]
\end{frame}


\begin{frame}{\textbf{Solution: Simplifying the Expression}}
We begin with:
\[
\sqrt[5]{70n \left( \sqrt[6]{70n} \right)^2}
= \sqrt[5]{70n \cdot (70n)^{2/6}}
= \sqrt[5]{(70n)^{1 + \frac{2}{6}}}
= \sqrt[5]{(70n)^{\frac{4}{3}}}
\]

Now apply the outer root:
\[
\left( (70n)^{\frac{4}{3}} \right)^{1/5} = (70n)^{\frac{4}{15}}
\]

We are told this is equal to:
\[
(70n)^{30x}
\Rightarrow 30x = \frac{4}{15}
\Rightarrow x = \frac{4}{450} = \frac{2}{225}
\]

\textbf{Answer:} \( \boxed{\frac{2}{225}} \)
\end{frame}

\begin{frame}
\section{Practices}
    \begin{tikzpicture}[remember picture, overlay]

        \shade[top color=softPurple, bottom color=softGray] 
            (current page.north west) rectangle (current page.south east);

        \fill[white, opacity=0.3] (current page.center) circle (4cm);
        

        \node[align=center] at (current page.center) {
            {\Huge\bfseries\textcolor{purpleblue}{Practices}} \\[0.5cm]
            {\Large\itshape\textcolor{lightText}{Excellence in SAT Prep}}
        };
        \foreach \x in {1,2,3,4,5} {
            \fill[purpleblue, opacity=0.2] (rand*10-5, rand*7-3.5) circle (0.1+\x*0.05);
        }
        ;
    \end{tikzpicture}
\end{frame}
\begin{frame}{Question}
If $\left(5^3\right)^{4k} = \left(5^{\frac{1}{3}}\right)^{24}$, what is the value of $k$?


\end{frame}

\begin{frame}{Answer}
We simplify both sides:

\[
(5^3)^{4k} = 5^{12k}, \quad \left(5^{\frac{1}{3}}\right)^{24} = 5^{\frac{24}{3}} = 5^8
\]

So,

\[
5^{12k} = 5^8 \Rightarrow 12k = 8 \Rightarrow k = \frac{2}{3}
\]

\textbf{Correct answer:} B
\end{frame}


\begin{frame}{Question}
Which expression is equivalent to $4^{m+2}$?

A) $16^m$ \\
B) $16 + 4^m$ \\
C) $8(4^m)$ \\
D) $16(4^m)$
\end{frame}


\begin{frame}{Answer}
We simplify the original expression:

\[
4^{m+2} = 4^m \cdot 4^2 = 4^m \cdot 16 = 16(4^m)
\]

\textbf{Correct answer:} D
\end{frame}

\begin{frame}{Question}
\[
\sqrt[8]{41k} \left( \sqrt[9]{41k} \right)^{12}
\]

For what value of $x$ is the given expression equivalent to $(41k)^{25x}$, where $k > 1$?
\end{frame}

\begin{frame}{Answer}


\textbf{Correct answer:} $x = \dfrac{7}{120}$
\end{frame}

\begin{frame}{Question}
\[
\frac{10 \sqrt[6]{9^3 \cdot x^{48}}}{\sqrt[4]{6^4 x}}
\]

The given expression is equivalent to $a x^b$, where $a > 0$, $b > 0$, and $x > 1$. What is the value of $ab$?
\end{frame}

\begin{frame}{Answer}
We simplify numerator and denominator:

Numerator:
\[
10 \cdot \sqrt[6]{9^3 \cdot x^{48}} = 10 \cdot (9^3 \cdot x^{48})^{1/6} = 10 \cdot 9^{1/2} \cdot x^8
\]

Since $9 = 3^2$, so $9^{1/2} = 3$, hence numerator becomes:
\[
10 \cdot 3 \cdot x^8 = 30x^8
\]

Denominator:
\[
\sqrt[4]{6^4 x} = (6^4 x)^{1/4} = 6 \cdot x^{1/4}
\]

Entire expression:
\[
\frac{30x^8}{6x^{1/4}} = 5x^{8 - 1/4} = 5x^{\frac{31}{4}}
\]

So $a = 5$, $b = \frac{31}{4}$

\[
ab = 5 \cdot \frac{31}{4} = \frac{155}{4}
\]


\end{frame}


\begin{frame}{Question}
\[
\sqrt[4]{x^3} \cdot \sqrt[12]{x^5}
\]

The given expression is equivalent to $\sqrt[n]{x^m}$, where $m$ and $n$ are positive constants each less than 10 and $x > 1$. What is the value of $\dfrac{m}{n}$?
\end{frame}


\begin{frame}{Answer}
We simplify the expression:

\[
\sqrt[4]{x^3} \cdot \sqrt[12]{x^5} = x^{\frac{3}{4}} \cdot x^{\frac{5}{12}} = x^{\frac{3}{4} + \frac{5}{12}}
\]

To add the exponents:

\[
\frac{3}{4} = \frac{9}{12}, \quad \frac{9}{12} + \frac{5}{12} = \frac{14}{12} = \frac{7}{6}
\]

So the expression becomes:
\[
x^{\frac{7}{6}} = \sqrt[6]{x^7}
\]

Therefore, $m = 7$, $n = 6$, and

\[
\frac{m}{n} = \frac{7}{6}
\]

\textbf{Correct answer:} $\dfrac{7}{6}$
\end{frame}


















\begin{frame}{}
    \begin{tikzpicture}[remember picture, overlay]
        % 背景渐变色：蓝紫到白
        \shade[top color=novelPurple!40, bottom color=white]
            (current page.north west) rectangle (current page.south east);

        % 中心柔光
        \fill[white, opacity=0.15] (current page.center) circle (5cm);

        % 柔和光点点缀
        \foreach \x in {1,2,3,4,5} {
            \fill[novelPurple, opacity=0.12] (rand*10-5, rand*7-3.5) circle (0.1+\x*0.05);
        }
    \end{tikzpicture}

    \centering
    \vspace{5em}

    {\Huge \textbf{\textcolor{novelPurple}{Thank You!}}} \\[1em]
    {\large \textcolor{gray}{We appreciate your time and focus.}} \\[3em]

    {\LARGE \textbf{\textcolor{novelPurple}{\textsf{Novel Prep}}}}

    \vfill
\end{frame}




\end{document}